% !PRINTED 184
\weberpage{184}

%<<CONT
and from this it is concluded that the coefficients $A_\nu, B_\nu, C_\nu, D_\nu$
\gesperrt{are integral rational functions of $\varkappa^2$ of degree $\nu$}, and
from \S\,54, 2, 3, it follows moreover that $A_\nu, B_\nu$ are, for even $\nu$,
expressed rationally by $\varkappa^2\varkappa'^2$, while for odd
$\nu$ they are equal to the product of $(\varkappa'^2-\varkappa^2)$ with a rational function
 of $\varkappa^2\varkappa'^2$.

From $B_\nu$ one finds $C_\nu$ by means of the formula:

\weq{10}{C_\nu(\varkappa^2) = \varkappa'^{2\nu}B_\nu\left(\frac{-\varkappa^2}{\varkappa'^2}\right)}

and $D_\nu$ by applying (6) and (8):

\weq{11}{D_\nu(\varkappa^2) = (-1)^\nu C_\nu(\varkappa'^2) = (-1)^\nu \varkappa^{2\nu}B_\nu\left(\frac{-\varkappa'^2}{\varkappa^2}\right).}

Since the coefficients $A_\nu(0), B_\nu(0), C_\nu(0), D_\nu(0)$ can easily be
determined from (5), the system of formulae developed here yields
the coefficients $A_\nu, B_\nu, C_\nu, D_\nu$ without
further calculation up to $\nu=3$ inclusive.

One finds:
\[
\begin{aligned}
&A_0=1, \quad A_1=0, \qquad A_2=-\frac{2}{3}(1-\varkappa^2\varkappa'^2),\\
&B_0=1, \quad B_1=\frac{1}{3}(\varkappa'^2-\varkappa^2), \quad B_2=\frac{1}{3}(1+2\varkappa^2\varkappa'^2),\\
&C_0=1, \quad C_1=\frac{1}{3}(1+\varkappa^2), \quad C_2=\frac{1}{3}(\varkappa'^4-2\varkappa^2),\\
&D_0=1, \quad D_1=-\frac{1}{3}(1+\varkappa'^2), \quad D_2=\frac{1}{3}(\varkappa^4-2\varkappa'^2).\\
&A_3=-\frac{8}{9}(\varkappa'^2-\varkappa^2)(2+\varkappa^2\varkappa'^2),\\
&B_3=\frac{1}{9}(\varkappa'^2-\varkappa^2)(5-2\varkappa^2\varkappa'^2),\\
&C_3=\frac{1}{9}(1+\varkappa^2)(5\varkappa'^4+2\varkappa^2),\\
&D_3=-\frac{1}{9}(1+\varkappa'^2)(5\varkappa^4+2\varkappa'^2).
\end{aligned}
\]

Further coefficients can be calculated in the same way, though
more laboriously, by developing the expressions for the
$\sigma$-functions in \S\,32 in powers of $q$ and using the higher
powers of $q$ as the lowest one has been used here. \gesperrt{Weierstrass}
makes use, for the recurrent calculation of the coefficients,
of certain partial differential equations, similar to the one that we
shall derive in the next section for the function $\sigma$.
